\begin{lemma}[Minimum bad pair: last path-index case]
Let \((\varphi,T)\) be a minimum bad pair with \(\rho(T)>2\).  The case
\(\rho(T)=|E(T)|\) is impossible.
\end{lemma}
\begin{proof}
Write \(T=(y_0,e_1,y_1,\ldots,e_n,y_n)\) and
\(q=\rho(T)=n\).  Unfolding \noderef{TashTreeMinimumBadPair} and
\noderef{TashTreeBadPair} gives the fixed \(k\)-triple, the Tashkinov
condition, non-elementariness, and the universal exclusion of every
smaller bad pair.  Since \(q>2\), all prefixes used below have positive
length.

The last edge \(e_n\) does not join \(y_{n-1}\) to \(y_n\).  Otherwise
\(n-1\) would satisfy \noderef{TashTreeSuffixPath}, whereas
\noderef{TashTreePathIndex} defines \(q=n\) to be the least suffix-path
start.  By the attachment field of \noderef{TreeSequence}, \(e_n\) joins
\(y_n\) to some earlier displayed vertex; the preceding observation shows
that this earlier endpoint belongs to \(V(T_{n-2})\).

We establish three successive normal forms.  At each stage we retain the
same symbols for the current coloring and its minimum bad pair, but we
justify minimum badness before invoking any prefix result for the new
coloring.

\emph{First normal form.}
The proper prefix \(T_{n-1}\), constructed by
\noderef{TreeSequencePrefix}, is elementary by
\noderef{TashTreeProperPrefixElementary}, while \(T\) is not.  Expanding
\noderef{ColoredVertexSetElementary}, some color \(\beta\) is missing at
\(y_n\) and at a vertex \(y_j\), \(j<n\); otherwise every repeated
missing color of \(T\) would already occur twice in the elementary
prefix.

By \noderef{TashTreeFourUnusedColors}, the proper prefix \(T_{n-2}\) has
four colors in \noderef{TreeFreeColors}.  By
\noderef{TreeMissingColors} each is missing at a displayed vertex
\(y_m\), \(m\le n-2\), and by \noderef{TreeUsedColors} none colors
\(e_2,\ldots,e_{n-2}\).  The two further exact identities
\(e_{n-1},e_n\) exclude at most two colors.  Hence at least two of the
four are unused on every colored tree edge; choose one called
\(\alpha\) which is distinct from \(\beta\), and let \(y_m\), \(m\le
n-2\), miss it.

Suppose first that \(j=n-1\).  The color \(\beta\) is unused on
\(e_2,\ldots,e_{n-1}\).  Indeed, if one of these edges had color
\(\beta\), \noderef{TashkinovTree} would give an earlier
missing-\(\beta\) witness, distinct from \(y_{n-1}\), contradicting
the elementariness of \(T_{n-1}\) just obtained from
\noderef{TashTreeProperPrefixElementary}.  It also does not color
\(e_n\), since that edge is incident with the \(\beta\)-missing vertex
\(y_n\).  Thus both \(\alpha,\beta\) are unused through \(y_{n-1}\), and
\noderef{TashTreeControlledInterchange} applies at \(y_m,y_{n-1}\).

That component does not contain \(y_n\).  It is a path with the two
distinct controlled ends \(y_m,y_{n-1}\).  The vertex \(y_n\) also
misses \(\beta\), so \noderef{MissingColorIncidenceEndpoint} gives it
degree at most one in the same two-color incidence graph.  If \(y_n\)
lay in the component, the spanning path supplied by
\noderef{TwoColorComponentSpanningPath} would have three distinct
degree-at-most-one vertex-nodes, impossible for a path.  Therefore the
swap fixes every edge incident with \(y_n\), and \(\beta\) remains
missing there.  The conclusion of
\noderef{TashTreeControlledInterchange} makes the new pair minimum bad
and moves the other missing occurrence of \(\beta\) to \(y_m\).  We have
the first normal form in this subcase.

Now suppose \(j\le n-2\).  If \(\beta\) is unused on the whole displayed
tree, the desired form already holds.  Otherwise let \(Q\) be the
component \noderef{TwoColorComponent} through \(y_n\) for
\(\alpha,\beta\).  Since \(y_n\) misses \(\beta\),
\noderef{MissingColorIncidenceEndpoint} makes it an end, and
\noderef{TwoColorComponentSpanningPath} gives a simple spanning path for
\(Q\), which we orient from \(y_n\).

Suppose \(Q\) meets \(V(T_{n-2})\).  Along that spanning path, let \(y_h\)
be the first vertex-node in this prefix and retain the segment from
\(y_h\) to \(y_n\).  Its internal vertex-nodes lie outside the retained
prefix by the choice of \(y_h\); simplicity gives distinct new vertices
and exact edge identities.  Append the segment, in reverse traversal from
\(y_h\) toward \(y_n\), to the tree-sequence \(T_{n-2}\).  The first
appended edge attaches its new vertex to the earlier \(y_h\), and every
later appended edge attaches the next new vertex to the preceding one, so
\noderef{TreeSequence} holds.  Every appended edge is colored
\(\alpha\) or \(\beta\), and the retained prefix contains the respective
missing witnesses \(y_m,y_j\).  Thus
\noderef{TashkinovTree} holds for the appended sequence.  Its unchanged
\(k\)-triple and the two missing-\(\beta\) vertices \(y_j,y_n\) make it
bad by \noderef{TashTreeBadPair}.  Its first appended edge may attach
away from the last retained vertex, but all subsequent appended edges are
in path order; therefore \(n-1\) is a suffix-path start and its path
index is at most \(n-1<q\).  By \noderef{TashTreeRankLess}, this is a
smaller bad pair, contrary to \noderef{TashTreeMinimumBadPair}.

It follows that \(Q\) avoids \(V(T_{n-2})\).  Every displayed tree edge
has an endpoint in that prefix.  For \(e_1,\ldots,e_{n-2}\) this follows
from both endpoints being prefix vertices; for \(e_{n-1}\), its earlier
endpoint is in the prefix; and for \(e_n\), it is the earlier-end fact
proved in the first paragraph.  If the edge-node of any displayed tree
edge belonged to \(Q\), its incidence edge in
\noderef{TwoColorIncidenceGraph} would put that prefix endpoint in \(Q\).
Thus no displayed exact tree-edge identity belongs to \(Q\).

Let \(\varphi_1\) be the component swap
\noderef{TwoColorComponentSwap} on \(Q\), and let \(\varphi_2\) be obtained
from \(\varphi_1\) by globally transposing the two color names
\(\alpha,\beta\), fixing all other colors and the uncolored value.  We
now verify, without invoking the controlled-interchange lemma, that
\((\varphi_2,T)\) is a minimum bad pair.

The local swap is proper by \noderef{TwoColorComponentSwapProper}; the
global transposition is an injective permutation of colored values, so
\noderef{IsProperEdgeColoring} remains true.  Both operations fix
\(\operatorname{none}\) and send colored values to colored values, hence
the same exact \(e_0\) remains uniquely uncolored by
\noderef{IsOneEdgeUncolored}.  All other clauses of \noderef{KTriple}
depend only on \(G,k,e_0\) and are unchanged.

For the Tashkinov condition under the local swap, first consider an edge
through \(T_{n-2}\).  Its exact identity and its old witness vertex both
avoid \(Q\), so its color and missing witness are fixed.  The only
remaining displayed edges are \(e_{n-1},e_n\), whose identities also
avoid \(Q\).  A witness for a color outside
\(\{\alpha,\beta\}\) is unchanged.  If either edge has color \(\alpha\)
or \(\beta\) and its old witness changes, use instead the fixed earlier
missing vertex \(y_m\) for \(\alpha\) or \(y_j\) for \(\beta\);
both lie in \(T_{n-2}\) and outside \(Q\).  Hence every edge retains an
earlier witness under \(\varphi_1\).  The global transposition then
changes an edge color and its missing-color relation by the same
permutation at every vertex, so the same witness position works under
\(\varphi_2\).  This proves \noderef{TashkinovTree}.

Under \(\varphi_1\), \(\alpha\) is missing at \(y_n\): if \(Q\) is
nontrivial, the unique incident \(\alpha\)-edge at the
\(\beta\)-missing end becomes \(\beta\); if \(Q\) is isolated, both
colors were already missing.  The vertex \(y_m\) lies outside \(Q\) and
continues to miss \(\alpha\).  After the global transposition,
\(\beta\) is missing at the two distinct vertices \(y_m,y_n\).
Thus \(V(T)\) is non-elementary and \noderef{TashTreeBadPair} holds for
\((\varphi_2,T)\).  Its sequence and rank are unchanged.  Finally, the
universal minimality conjunct of the original
\noderef{TashTreeMinimumBadPair} quantifies over every coloring and
therefore combines with this badness to make \((\varphi_2,T)\) minimum.

No tree-edge identity was changed by the local swap.  Consequently, if a
tree edge has color \(\beta\) under \(\varphi_2\), it had color
\(\alpha\) under the original coloring, contrary to the choice of
\(\alpha\).  We have therefore reached the first normal form:
\[
 \beta\in\overline\varphi(y_j)\cap\overline\varphi(y_n)
 \quad\text{for some }j\le n-2,\qquad
 \beta\notin\varphi\langle T\rangle .
\]

\emph{Second normal form.}
Put \(\gamma=\varphi(e_n)\).  Suppose it is missing at \(y_{n-1}\).
The Tashkinov condition for \(e_n\) supplies a missing-\(\gamma\)
witness in \(T_{n-1}\).  Since that prefix is elementary by
\noderef{TashTreeProperPrefixElementary} for the current minimum bad
pair, the witness must be \(y_{n-1}\) itself.  Moreover \(\gamma\) is
unused on \(e_2,\ldots,e_{n-1}\): any occurrence there would have an
earlier missing-\(\gamma\) witness by \noderef{TashkinovTree}, distinct
from \(y_{n-1}\), and would contradict the same prefix elementariness.

The first-normal-form color \(\beta\) is unused on all tree edges.
Apply \noderef{TashTreeControlledInterchange} to \(\beta,\gamma\) at
the missing vertices \(y_j,y_{n-1}\).  Its controlled component cannot
contain \(y_n\).  Indeed, \(y_n\) misses \(\beta\) and is incident with
the \(\gamma\)-colored exact edge \(e_n\); by
\noderef{MissingColorIncidenceEndpoint} it has degree at most one, and
the displayed incidence makes that degree one.  The spanning path from
\noderef{TwoColorComponentSpanningPath} already has the two distinct
ends \(y_j,y_{n-1}\), so \(y_n\) cannot be a third end.  Thus \(y_n\)
and \(e_n\) are outside the component, and \(e_n\) retains color
\(\gamma\).  At the other controlled end \(y_{n-1}\), the nontrivial
path has an incident \(\beta\)-edge, which becomes \(\gamma\).  Hence
\(\gamma\) is present at \(y_{n-1}\).  The controlled lemma also supplies
minimum badness of the new pair.  This proves the second normal form.

\emph{Third normal form.}
Using non-elementariness of the current \(T\) and elementariness of its
proper prefix from \noderef{TashTreeProperPrefixElementary}, choose a
color \(\beta'\) missing at \(y_n\) and at \(y_{j'}\), \(j'<n\).
If \(j'\le n-2\), the third form already holds.  If \(j'=n-1\), the same
prefix-elementarity argument is as follows.  If some
\(e_h\), \(2\le h\le n-1\), had color \(\beta'\), then
\noderef{TashkinovTree} would supply a missing-\(\beta'\) witness
\(y_r\) with \(r<h\le n-1\).  This vertex is distinct from
\(y_{n-1}\), contradicting the elementariness of \(T_{n-1}\) supplied by
\noderef{TashTreeProperPrefixElementary}.  Thus \(\beta'\) is unused on
\(e_2,\ldots,e_{n-1}\), and incidence with its missing vertex \(y_n\)
shows that it is not the color of \(e_n\).

Apply \noderef{TashTreeFourUnusedColors} to \(T_{n-2}\) for the current
minimum bad pair.  After excluding the at most two colors on
\(e_{n-1},e_n\), at least two free colors remain unused on the whole
tree.  Choose one \(\alpha'\ne\beta'\), missing at a vertex of
\(T_{n-2}\).  A controlled \((\alpha',\beta')\)-interchange at that
vertex and \(y_{n-1}\) moves missing \(\beta'\) into \(T_{n-2}\).
Exactly as in the first controlled swap, \(y_n\), which misses
\(\beta'\), is a third degree-at-most-one vertex and hence lies outside
the component by \noderef{MissingColorIncidenceEndpoint} and
\noderef{TwoColorComponentSpanningPath}; thus it continues to miss
\(\beta'\).

Neither affected color is \(\gamma\).  The color \(\alpha'\) is unused
on \(e_n\), which has color \(\gamma\); and \(\beta'\) is missing at
\(y_n\), where the incident \(e_n\) makes \(\gamma\) present.  Therefore
the interchange neither changes \(e_n\)'s color nor the presence of
\(\gamma\) at \(y_{n-1}\).  The second normal form survives, and the
conclusion of \noderef{TashTreeControlledInterchange} again supplies a
minimum bad pair.  The third normal form is established.

Finally delete the exact pair \(e_{n-1},y_{n-1}\), retaining
\[
 T'=(y_0,e_1,y_1,\ldots,e_{n-2},y_{n-2},e_n,y_n).
\]
This is a tree-sequence: deleting entries preserves vertex and edge
injectivity and the first edge, all earlier attachment equations are
unchanged, and the earlier endpoint of \(e_n\) lies in \(T_{n-2}\).
All Tashkinov witnesses before the final edge are inherited.  The final
edge still has color \(\gamma\).  Its old missing-\(\gamma\) witness
cannot be \(y_{n-1}\), because the second normal form says \(\gamma\) is
present there, and \noderef{PresentColors} and \noderef{MissingColors}
are complementary.  Hence a witness for \(e_n\) remains in
\(T_{n-2}\), proving \noderef{TashkinovTree} for \(T'\).

The \(k\)-triple is unchanged, and the third normal form supplies a color
missing both at \(y_n\) and at a retained vertex of \(T_{n-2}\).  Thus
\noderef{TashTreeBadPair} holds for the shortened pair.  It has \(n-1\)
edges, so \noderef{TashTreePathIndex} gives
\(\rho(T')\le n-1<n=q\).  By \noderef{TashTreeRankLess}, it is smaller
than the current minimum bad pair, contradicting
\noderef{TashTreeMinimumBadPair}.  Therefore the case \(q=n\) is
impossible.
\end{proof}
